% 原文 pp.53–58（印刷 pp.1171–1176）
\section{Lecture 4：二维 $\sigma$ 模型的红外行为与 $S$ 矩阵}

\begin{keytakeaways}
本讲研究靶空间为球面 $S^{N-1}$ 的二维 $\sigma$ 模型的红外行为。物理图像是：该模型在二维可重整化、且映入正曲率空间时 $\beta$ 函数为负（紫外渐近自由），故经典边缘耦合在红外变为相关，微扰论所预期的 $N-1$ 个无质量 Goldstone 玻色子并不出现；正确答案是 $N$ 个质量同为 $M$ 的有质量粒子，红外极限平庸。核心数学工具是：用 $\delta$ 函数的 Fourier 表示引入辅助标量场 $\sigma$，把对球面（弯曲靶空间）的积分化为对线性空间的 Gauss 积分，再取 $N\to\infty$ 的鞍点近似；由此得到谱条件 $\frac{\tilde\lambda}{4\pi}\ln\frac{\Lambda^2}{-i\tilde\lambda\rho}=1$、量纲透射给出的质量 $M=\Lambda e^{-2\pi/\tilde\lambda}$，以及 $\beta$ 函数只有单圈贡献这一结论。本讲后半转入 $S$ 矩阵：二维时 Coleman–Mandula 定理失效，存在「被禁戒的」对称性；作者用应力-能量张量的双分次（滤子）结构证明守恒方程 $\partial_-T_{++}+\partial_+T_{+-}=0$ 可以量子地成立，并由此出发构造非平凡的守恒流，进而显式求出 $S$ 矩阵。
\end{keytakeaways}

% ===== 原文 p.53 =====
\subsection{某些二维 $\sigma$ 模型的红外行为}
% 原文 §4.1：Infrared behavior of certain two-dimensional sigma models

\noindent\textit{（本讲原标题：靶空间为 $S^{N-1}$ 的二维 $\sigma$ 模型的红外行为与 $S$ 矩阵。）}

% 原文标题：Infrared Behavior and the S-matrix of the Two-Dimensional Sigma Model with Target Space $S^{N-1}$

考虑 $N$ 个标量玻色子的理论，其拉格朗日量为
\begin{equation}
  \tag{4.1}
  \label{eq:4-1}
  \mathcal{L}=\int d^dx\left(\frac{1}{2}(\nabla\phi)^2+\frac{g}{4!}(\phi^2-a^2)^2\right),
\end{equation}
其中 $\phi:\mathbb{R}^d\to\mathbb{R}^N$。在第 1 讲中我们看到：当 $d>2$、且 $g$ 充分小时，这个理论具有对称性破缺（spontaneous symmetry breaking），对称群从 $SO(N)$ 破缺到 $SO(N-1)$。特别地，对球面 $\phi^2=a^2$ 上的任一点 $y$，都存在该理论的一个实现（realization）$\mathcal{H}_y$，使得 $\langle\phi\rangle=y$。我们还看到，(2.1) 的低能有效理论（effective theory）是靶空间（target space）为 $S^{N-1}$ 的 $\sigma$ 模型（sigma model）。

在 2 维情形，第一个论断失效：不存在对称性破缺，并且只有一个实现，它带有 $SO(N)$ 的作用。然而第二个论断仍然成立：低能有效理论就是 $\sigma$ 模型。

在 $d>2$ 时我们知道 $\sigma$ 模型是红外自由（infrared free）的（在弱耦合近似下）；更确切地说，它在红外极限下收敛为 $N-1$ 个无质量标量场（即 Goldstone 玻色子）的自由理论。2 维情形并非如此。的确，在 2 维中 $\sigma$ 模型是可重整化的，因此它具有经典意义下边缘（marginal）的相互作用（耦合是无量纲的）。另一方面，[II-CFT] 第 3 讲已经证明：映入正曲率空间的二维 $\sigma$ 模型的 $\beta$ 函数为负（即该模型在紫外极限下渐近自由（asymptotic freedom））。因此，这些经典边缘的相互作用在红外是相关的（relevant），我们无法断言该模型红外自由。事实上，实际发生的情形（今天就会看到）是：与 $d>2$ 时拥有 $N-1$ 个无质量粒子不同，我们将拥有 $N$ 个有质量粒子，于是红外极限是平庸的（关联函数在动量空间原点解析，在位置空间无穷远处指数衰减），且与微扰论所见颇为不同。

在本讲前半部分，我们将证明二维 $\sigma$ 模型在大 $N$ 时的确红外平庸。我们还将建立一套以 $1/N$ 为级数展开计算关联函数的程序。

\noindent\textit{（本讲笔记由 Pavel Etingof 与 David Kazhdan 整理。）}

% [Tutor Note: 此处原文写「(2.1) 的低能有效理论」，但本讲只出现过 (4.1) 这一个四次拉格朗日量，Lecture 4 全文中并无 (2.1)；对照第 1 讲的同一论断（那里确有一编号较早的拉格朗日量），可判定此处是排印错误，应读作「(4.1) 的低能有效理论」，即同一 $O(N)$ 线性 $\sigma$ 模型的低能极限。]

% [Tutor Note: 「在 2 维不存在对称性破缺」并非平凡陈述，其依据是 $d=1+1$ 维的 Mermin–Wagner–Coleman 定理：连续整体对称性在二维不能自发破缺（Sidney Coleman, Comm. Math. Phys. 31 (1973) 259）。证明的关键正是红外行为——红外涨落（二维的 $\int d^2k/k^2$ 型对数发散）把序参量的期望值抹平；因此这里「第一个论断失效、第二个论断仍成立」并不是逻辑上的偶然，而是 $d=2$ 与 $d\ge3$ 的红外结构本质不同所致。]

% ===== 原文 p.54 =====
\textbf{注 1（Remark 1）.} 由于映入球面的二维 $\sigma$ 模型的 $\beta$ 函数为负，人们相信这个理论确实存在。能够在 $1/N$ 展开中构造该理论——它给出的信息与标准微扰论所见相当不同，却又与关于重整化群的标准理论相符——为这一信念提供了强有力的支持。

\textbf{注 2（Remark 2）.} 由于四次拉格朗日量 (4.1) 所定义的理论在 $d=2$ 中严格存在，它可以被视为 $\sigma$ 模型的一个紫外截断（UV cutoff）。该截断的特征动量标度 $\Lambda$ 就是 $\phi$ 的径向分量的质量。这个截断模型称为线性 $\sigma$ 模型（linear sigma model），以区别于非线性 $\sigma$ 模型（nonlinear sigma model）——后者是映入球面的映射所构成的模型。

\subsection{$N\to\infty$ 极限下红外行为的计算}
% 原文 §4.2：Computation of the infrared behavior in the $N\to\infty$ limit

我们将这样来计算 $\sigma$ 模型在 $N\to\infty$ 极限下的红外行为：先把积分整理成 $N$ 只作为参数出现的形式，再对路径积分作鞍点近似（saddle point approximation）。我们将形式地处理路径积分，不去操心它们是否存在。

定义 $\sigma$ 模型欧氏关联函数生成级数的路径积分由下式给出：
\begin{equation}
  \tag{4.2}
  \label{eq:4-2}
  Z(J)=\int_{\phi:\,\phi^2=1}D\phi\,\exp\left(-\frac{1}{2\lambda}\int d^2x\,|d\phi|^2+Z_J\int d^2x\,J\phi\right),
\end{equation}
其中 $\phi,J:\mathbb{R}^2\to\mathbb{R}^N$，$J\phi$ 是 $(J,\phi)$ 的简写，$\lambda$ 是耦合常数。加入因子 $Z_J$ 是因为算符 $\phi$ 具有反常维数（anomalous dimension，见 [I-Renorm Groups] 第 3 讲），需要作乘性重整化。（如我们所知，$Z_J$ 依赖于截断，并在截断趋于无穷时发散。）

对映入线性空间的映射作积分，比对映入弯曲流形的映射作积分更方便。因此，把积分 (4.2) 改写成下述形式是有用的：
\begin{equation}
  \tag{4.3}
  \label{eq:4-3}
  Z(J)=\int_{\phi:\,\mathbb{R}^2\to\mathbb{R}^N}D\phi\prod_{x\in\mathbb{R}^2}\delta\bigl(\phi^2(x)-1\bigr)\exp\left(-\frac{1}{2\lambda}\int d^2x\,|d\phi|^2+Z_J\int d^2x\,J\phi\right).
\end{equation}

现在我们来用公式
\begin{equation}
  \tag{4.4}
  \label{eq:4-4}
  \delta(a)=\frac{1}{2\pi}\int_{-\infty}^{\infty}e^{ia\sigma}\,d\sigma .
\end{equation}
把 (4.4) 代入 (4.3)，得到（差一个与 $J$ 无关的因子）：
\begin{equation}
  \tag{4.5}
  \label{eq:4-5}
  Z(J)=\int D\phi\int D\sigma\,\exp\left(i\int d^2x\,\frac{\sigma}{2}\bigl(\phi^2-1\bigr)-\frac{1}{2\lambda}\int d^2x\,|d\phi|^2+Z_J\int d^2x\,J\phi\right).
\end{equation}
于是，使用公式 (4.4) 引入了辅助标量场 $\sigma$。

现在我们将交换积分次序，先对 $\phi$ 作积分。我们的目标是把积分约化为 $N$（靶空间的维数）只作为参数进入的形式。幸运的是，球面是由一个二次方程定义的，因此在刚才写出的形式中，对 $\phi$ 的积分是高斯型的。

% [Tutor Note: (4.4) 给出的是指数 $i\sigma(\phi^2-1)$，而 (4.5) 中出现的却是 $\frac{\sigma}{2}(\phi^2-1)$。这里并非漏抄因子 $2$：作者把 (4.4) 中的积分变量 $\sigma$ 重命名为 $\sigma/2$（常数 Jacobian 被并入「与 $J$ 无关的因子」）。这一重命名正是后面 (4.6) 中出现 $-\frac{\Delta}{\lambda}-i\sigma$ 这一组合的原因：把 $\frac{i\sigma}{2}\phi^2$ 与 $-\frac{1}{2\lambda}|d\phi|^2$ 合并，得 $\frac{1}{2\lambda}\phi(-\Delta-i\lambda\sigma)\phi$，于是算符的「质量项」写成 $i\lambda\sigma$ 而非 $2i\lambda\sigma$。]

% ===== 原文 p.55 =====
因此我们可以显式地算出这个积分，得到
\begin{equation}
  \tag{4.6}
  \label{eq:4-6}
  Z(J)=\int D\sigma\,\det\left(-\frac{\Delta}{\lambda}-i\sigma\right)^{-N/2}\exp\left(-i\int d^2x\,\frac{\sigma}{2}-\frac{1}{2}Z_J^2\int d^2x\,J\cdot\left(\frac{\Delta}{\lambda}+i\sigma\right)^{-1}J\right)
\end{equation}
（行列式的幂次 $-N/2$ 之所以出现，是因为我们对 $\phi$ 的各个分量 $\phi_i$ 作了 $N$ 个独立的高斯积分）。把行列式提到指数上去是有用的：
\begin{align}
  \tag{4.7}
  \label{eq:4-7}
  Z(J)&=\int D\sigma\,\exp\Bigl(-\frac{N}{2}\operatorname{Tr}\ln\Bigl(-\frac{\Delta}{\lambda}-i\sigma\Bigr)\Bigr.\nonumber\\
  &\qquad\qquad\Bigl.-i\int d^2x\,\frac{\sigma}{2}-\frac{1}{2}Z_J^2\int d^2x\,J\cdot\Bigl(\frac{\Delta}{\lambda}+i\sigma\Bigr)^{-1}J\Bigr).
\end{align}

现在我们要把积分 (4.6) 化成这样的形式：指数中每一项前面都带因子 $N$（$J$ 项除外——它用于把答案按 $J$ 的幂作形式级数展开；当 $N\to\infty$ 时我们令 $J$ 为 $O(1)$，使得这一项对大 $N$ 而言是小微扰，不影响鞍点）。然后我们用鞍点近似来计算该积分。为此我们必须作变量替换 $\rho=\sigma/N$，并把耦合常数重新标度为 $\tilde\lambda=\lambda N$。经过这些改变，积分（差一个因子）形如
\begin{align}
  \tag{4.8}
  \label{eq:4-8}
  Z(J)&=\int D\rho\,\exp\Bigl(-\frac{N}{2}\Bigl[\operatorname{Tr}\ln\bigl(-\Delta-i\tilde\lambda\rho\bigr)+i\int d^2x\,\rho\Bigr]\Bigr.\nonumber\\
  &\qquad\qquad\Bigl.-\frac{\tilde\lambda}{2N}Z_J^2\int d^2x\,J\cdot\bigl(\Delta+i\tilde\lambda\rho\bigr)^{-1}J\Bigr).
\end{align}

在这个形式中，我们已经可以应用鞍点近似了。（这里确实是鞍点而非驻定相位（stationary phase），因为指数中的函数是复值的。）

寻找庞加莱不变的鞍点是自然的，因为我们所考虑的积分是庞加莱不变的。因此，我们应当寻找一个为常数的鞍点 $\rho(x)$，即 $\rho(x)=\rho$。$\rho$ 的值由鞍点方程
\begin{equation}
  \tag{4.9}
  \label{eq:4-9}
  \frac{d}{d\rho}\left[\operatorname{Tr}\ln\bigl(-\Delta-i\tilde\lambda\rho\bigr)+i\rho\int d^2x\right]=0
\end{equation}
定出。（此时我们临时引入一个红外截断 $L$，即假定时空是尺度为 $L\times L$ 的环面。这使总作用量有限。）让我们用算符 $-\Delta-i\tilde\lambda\rho$ 的动量空间表示来改写 (4.9)，并令 $L\to\infty$。于是得到
\begin{equation}
  \tag{4.10}
  \label{eq:4-10}
  \tilde\lambda\int\frac{d^2k}{(2\pi)^2}\,\frac{1}{k^2-i\tilde\lambda\rho}=1 .
\end{equation}
左端的积分是对数型紫外发散的。因此我们应在动量空间中引入紫外截断 $\Lambda$（即只对球 $|k|<\Lambda$ 作积分）。于是 (4.10) 变为
\begin{equation}
  \tag{4.11}
  \label{eq:4-11}
  \frac{\tilde\lambda}{4\pi}\ln\frac{\Lambda^2}{-i\tilde\lambda\rho}=1 .
\end{equation}

% [Tutor Note: 从 (4.5) 到 (4.6) 的符号是本节最易出错之处。对 $\phi$ 的高斯积分中，二次型算符是 $A'\equiv\frac{1}{\lambda}(-\Delta-i\lambda\sigma)$，而公式给出 $\det(A')^{-1/2}\exp\bigl(+\frac12 b\,A'^{-1}b\bigr)$；由于 $(\Delta/\lambda+i\sigma)^{-1}=-\bigl(-\Delta/\lambda-i\sigma\bigr)^{-1}$ 自带一个负号，它恰与 $A'^{-1}=\lambda(-\Delta-i\lambda\sigma)^{-1}$ 的符号结构相合，才得到 (4.6) 中整体为负的 $-\frac12Z_J^2\int J\cdot(\Delta/\lambda+i\sigma)^{-1}J$。若改用闵可夫斯基号差或把 $\exp(+J\phi)$ 的符号写反，这一项会整体变号，进而使 (4.10) 右端变成 $-1$，谱条件随之改写。]

% [Tutor Note: (4.8) 中「让 $N$ 只作参数」的两步替换 $\rho=\sigma/N$、$\tilde\lambda=\lambda N$ 有一个容易被忽略的副产品：$-\frac{\Delta}{\lambda}-i\sigma=\frac{N}{\tilde\lambda}\bigl(-\Delta-i\tilde\lambda\rho\bigr)$，取 $\operatorname{Tr}\ln$ 时会产生一个与场无关的常数 $\frac{N}{2}\operatorname{Tr}\ln\frac{N}{\tilde\lambda}$（即 $\frac{N}{2}\int d^2x\int\frac{d^2k}{(2\pi)^2}\ln\frac{N}{\tilde\lambda}$）。作者把它归入「差一个因子」，但严格说来它是发散的、且依赖截断；只是在鞍点方程 (4.9) 中对 $\rho$ 求导时它消失，故不影响谱条件。]

% [Tutor Note: (4.9) 到 (4.10) 用了 $\frac{d}{d\rho}\operatorname{Tr}\ln(-\Delta-i\tilde\lambda\rho)=-\operatorname{Tr}\bigl[(-i\tilde\lambda)(-\Delta-i\tilde\lambda\rho)^{-1}\bigr]$，其中泛函迹在环面上给出 $\operatorname{Tr}=\int d^2x\int\frac{d^2k}{(2\pi)^2}$，时空体积 $L^2$ 与 $i\rho\int d^2x$ 的体积因子相消后才有 (4.10)——这一步若忘记体积因子，会误以为谱条件依赖于环面大小 $L$。另外 (4.11) 的系数 $\frac{1}{4\pi}$ 来自 $\int_{|k|<\Lambda}\frac{d^2k}{(2\pi)^2}\frac{1}{k^2+m^2}=\frac{1}{4\pi}\ln\frac{\Lambda^2+m^2}{m^2}\approx\frac{1}{4\pi}\ln\frac{\Lambda^2}{m^2}$（角向积分给 $2\pi$，径向给 $\frac{1}{2}$），少乘 $2\pi$ 或漏掉径向的 $\frac12$ 都会把 $\frac{1}{4\pi}$ 写成 $\frac{1}{2\pi}$ 或 $\frac{1}{8\pi^2}$。]

% ===== 原文 p.56 =====
这表明
\begin{equation}
  \tag{4.12}
  \label{eq:4-12}
  \rho=\frac{i\Lambda^2}{\tilde\lambda}\exp\left(-\frac{4\pi}{\tilde\lambda}\right).
\end{equation}

在 [I-Renorm Groups] 第 3 讲中已经证明：在一个渐近自由的理论中，要使物理在 $\Lambda\to\infty$ 时有极限，耦合必须按下式调节：
\begin{equation}
  \tag{4.13}
  \label{eq:4-13}
  \tilde\lambda(\Lambda)=\frac{1}{A\ln\Lambda}+o\Bigl(\frac{1}{\ln\Lambda}\Bigr),\qquad \Lambda\to\infty,
\end{equation}
其中 $A$ 是 $\tilde\lambda$ 的单圈 $\beta$ 函数系数的负值。在 $\sigma$ 模型中可以算出 $A=(2\pi)^{-1}\frac{N-2}{N}$，因此它模 $1/N$ 等于 $(2\pi)^{-1}$（见 [II-CFT] 第 3 讲末尾）。我们刚刚导出的 $\rho$ 的公式表明：模 $1/N$，$\beta$ 函数等于它的单圈近似。因为，当且仅当 $\beta$ 函数模 $1/N$ 取此值时，由 (4.12) 给出的 $\rho$ 才满足重整化群方程——它当然应当满足。\footnote{也可以更直接地通过 Feynman 图证明：在大 $N$ 极限下，$\beta$ 函数只有单圈贡献。证明这一点只需画图——不需要任何计算。把它想清楚，是推荐大家做的一个练习。}重整化是这样实现的：令 $\Lambda\to\infty$、$\tilde\lambda\to0$，同时保持重整化群不变量 $M^2=\lim_{\Lambda\to\infty}\bigl(-i\tilde\lambda\rho\bigr)$ 固定；$M$ 最后将成为粒子质量。

\textbf{注（Remark）.} 于是，在量子理论中我们有了一个由 $M$ 定义的特征动量标度，尽管经典理论是共形不变的，且拉格朗日量只含一个无量纲耦合 $\lambda$、并未定义任何特定的动量（或长度）标度。这一现象称为「量纲透射」（dimensional transmutation）。它非常生动地说明：对一个具有标度不变拉格朗日量的理论作量子化，一般并不保持标度不变性，因而量子化并不完全是一件标准（canonical）的事情。

回想 $Z(J)$ 中 $J$ 项的形式，我们得出结论：$Z(J)$ 展开中的 $J^2$ 项是 $1/N$ 阶的。因此，为了得到 $N\to\infty$ 时关联函数的非零极限，我们应重新标度场，即 $\phi\to N^{1/2}\phi$。作了这一改变之后可以算出：$\phi_i$ 的两点函数有非零极限，而连通的 4、6、$\dots$ 点函数模 $1/N$ 为零。于是，在 $1/N$ 的一阶近似下，我们的理论是自由的。而且，这一阶的传播子是 $\frac{1}{k^2+M^2}$，其中 $M$ 如上定义。因此，极限理论是 $N$ 个独立的、质量为 $M$ 的自由有质量标量场的理论。

连通四点函数是 $1/N$ 阶的。因此，有质量标量之间的有效耦合常数是 $1/N$ 阶的。

场 $\phi_i$ 在大 $N$ 极限下退耦并变成有质量标量场，这一事实有如下概率论解释。

考虑 $N$ 个自由有质量玻色子 $a_i$ 的理论，它由路径积分
\begin{equation}
  \tag{4.14}
  \label{eq:4-14}
  Z=\int\Bigl(\prod Da_i\Bigr)\exp\Bigl(-N\sum_i\bigl(|da_i|^2+M^2a_i^2\bigr)\Bigr).
\end{equation}
描述。

% [Tutor Note: (4.12) 的 $\rho$ 是纯虚数（$\rho=i\Lambda^2\tilde\lambda^{-1}e^{-4\pi/\tilde\lambda}$），这正是作者在 (4.8) 后强调「这是鞍点而非驻定相位」的原因：指数中的函数为复值，鞍点不在实轴上，$\rho$ 的虚部使 $i\tilde\lambda\rho=-M^2$ 恰为实数且为负，从而 $-i\tilde\lambda\rho=M^2>0$ 起「质量平方」的作用。若误把 $\rho$ 当实数去求极小，会得到完全错误的符号结构。]

% [Tutor Note: 作者从 (4.11) 直接给出 $A=(2\pi)^{-1}$，未写出检验，但检验很短：由 (4.12) 有 $M^2=-i\tilde\lambda\rho=\Lambda^2e^{-4\pi/\tilde\lambda}$，要求 $\Lambda\to\infty$ 时 $M$ 固定，就必须 $\frac{4\pi}{\tilde\lambda}\approx 2\ln\Lambda$，即 $\tilde\lambda\approx\frac{2\pi}{\ln\Lambda}$，故 $\tilde\lambda=\frac{1}{A\ln\Lambda}$ 中 $A=\frac{1}{2\pi}$；这与单圈 $\beta$ 函数 $\frac{d\tilde\lambda}{d\ln\Lambda}=-A\tilde\lambda^2$ 的系数一致。也正由此看出 $\frac{N-2}{N}$ 的修正属于 $1/N$ 阶，与「$\beta$ 函数模 $1/N$ 等于单圈」的说法相符。]

% [Tutor Note: (4.14) 的指数中 $\sum_i$ 之后没有 $\int d^2x$：由于 $|da_i|^2$ 是二次微分形式，作用量必须对时空积分，故此处的 $\sum_i$ 应读作 $\sum_i\int d^2x$ 的简写（与 (4.15) 中传播子 $\frac{1}{N(k^2+M^2)}$ 的归一化一致：指数里整体带因子 $N$，而每个 $a_i$ 的传播子因此是 $\frac{1}{N(k^2+M^2)}$）。这是原文的书写疏漏，不是新的约定。]

% ===== 原文 p.57 =====
考虑该理论中算符 $a_i^2$（固定某个 $i$）的期望值 $\langle a_i^2\rangle$。这个一点函数需要重整化：它由积分
\begin{equation}
  \tag{4.15}
  \label{eq:4-15}
  \int\frac{d^2k}{(2\pi)^2}\,\frac{1}{N(k^2+M^2)}
\end{equation}
给出，而该积分紫外发散。因此我们应引入紫外截断 $\Lambda$。于是 $\langle a_i^2\rangle=\frac{1}{4\pi N}\ln\Lambda^2/M^2$。因此 $\bigl\langle\sum_i a_i^2\bigr\rangle=\frac{1}{4\pi}\ln\Lambda^2/M^2$。另一方面，由大数定律，若 $d$ 是 $a_i^2$ 的弥散（dispersion），则 $\sum a_i^2$ 的弥散为 $d/\sqrt{N}$，故对大 $N$，向量 $(a_1,\dots,a_N)$ 始终靠近球面 $\sum_i x_i^2=\frac{1}{4\pi}\ln\Lambda^2/M^2$。于是，$N$ 个独立有质量场的理论在大 $N$ 时的行为就像球面的 $\sigma$ 模型，其半径随截断作对数式变化。反过来，球面的 $\sigma$ 模型就是 $N$ 个有质量玻色子的自由理论，模 $1/N$ 阶的修正。

% [Tutor Note: 这一段有两处原文笔误，均不影响结论：(i)「stays close to the the sphere」重了一个 the；(ii) 球面方程写作 $\sum_i x_i^2$，而本段讨论的场是 $a_i$，按上下文应为 $\sum_i a_i^2$。另外「$d$ 是 $a_i^2$ 的弥散、$\sum a_i^2$ 的弥散为 $d/\sqrt N$」是把 $N$ 个独立同分布变量求和的涨落按 $\sqrt N$ 估计（中心极限定理），这正是「大 $N$ 时向量贴近半径为 $\frac{1}{4\pi}\ln\frac{\Lambda^2}{M^2}$ 的球面」的定量含义。]

\subsection{$S$ 矩阵的计算}
% 原文 §4.3：Computation of the S-matrix

现在我们将以完全不同的方式来研究映入球面的 $\sigma$ 模型。也就是说，我们将证明这个模型即便对有限的 $N$ 也是可积的，其含义是它的 $S$ 矩阵可以显式地算出。

回想 Coleman–Mandula 定理：在维数 $d>2$、有质量间隙且 $S$ 矩阵非退化的场论中，$S$ 矩阵的任一（哪怕是无穷小的）对称性，都是庞加莱李代数的一个元素与一个同庞加莱李代数对易的对称性的线性组合。换言之，若存在一个「被禁戒的」对称性（即不具有这种分解的对称性），则 $S$ 矩阵在某些技术条件下等于 $1$，亦即理论是自由的。

在二维情形，Coleman–Mandula 定理不成立。也就是说，确实存在 $S\ne1$ 的二维量子场论，它们具有一个不具有 Coleman–Mandula 定理所描述的那种分解的无穷小对称性。事实上，我们将看到靶空间为球面的二维 $\sigma$ 模型就是一例。

因此，在二维情形，「存在被禁戒的对称性」这一信息不足以断定 $S=1$。然而，这一信息仍然允许我们计算 $S$ 矩阵。本节我们将说明如何对映入 $S^{N-1}$ 的 $\sigma$ 模型做到这一点。

我们将在坐标为 $x_+,x_-$、度规为 $dx_+dx_-$ 的平坦闵可夫斯基时空中工作。令
\begin{equation}
  \tag{4.16}
  \label{eq:4-16}
  \mathcal{L}=\frac{1}{2\lambda}\int d^2x\,(\partial_+\phi\,\partial_-\phi),\qquad \phi^2=1
\end{equation}
为靶空间是 $S^{N-1}$（$\phi\in\mathbb{R}^N$）的 $\sigma$ 模型的拉格朗日量。由于这个拉格朗日量可以自然地写在弯曲时空上，它具有应力-能量张量（stress-energy tensor）$T=\frac{\delta\mathcal{L}}{\delta g}$。这是一个对称的二阶张量，故在坐标 $x_+,x_-$ 下有 3 个分量：$T_{++}$、$T_{--}$、$T_{+-}$。

让我们考察应力-能量张量的一些性质。经典地，我们有
\begin{equation}
  \tag{4.17}
  \label{eq:4-17}
  T_{++}=(\partial_+\phi)^2,\qquad T_{--}=(\partial_-\phi)^2 .
\end{equation}
（像 $(\partial_+\phi)^2$ 这样的表达式是 $(\partial_+\phi,\partial_+\phi)$ 的简写，其中的内积由 $\mathbb{R}^N$ 上的度规诱导。）由于 $\sigma$ 模型经典地共形不变，故 $T_{+-}=0$。

% [Tutor Note: (4.16)–(4.17) 采用光锥坐标 $x_\pm$ 与度规 $ds^2=dx_+dx_-$（即 $g_{+-}=g_{-+}=\frac12$，号差 $(+,-)$），而非许多教材的 $ds^2=2dx_+dx_-$；正是这个「不带因子 2」的约定，使得 $T_{++}$ 可以不加任何系数就等于 $(\partial_+\phi)^2$。若改用 $ds^2=2dx_+dx_-$，$T_{++}$ 会整体差一个因子 $2$，(4.16) 中的 $\frac{1}{2\lambda}$ 与欧氏 (4.2) 中 $-\frac{1}{2\lambda}\int|d\phi|^2$ 的对应关系也随之改变。另需注意 (4.2) 与 (4.16) 之间已作 Wick 转动：欧氏作用量里的 $|d\phi|^2$ 正定，而闵氏 (4.16) 的 $\partial_+\phi\partial_-\phi$ 含不定号差，两式不能逐项对照符号。]

% ===== 原文 p.58 =====
然而，我们应当预期：量子力学上共形不变性被破坏（事实上，这一点从大 $N$ 极限已经知道），所以在量子力学上 $T_{+-}\ne0$；事实上我们将看到，$T_{+-}$ 正比于 $\partial_+\phi\,\partial_-\phi$（见下文）。尽管如此，只要算符 $T_{+-}$ 定义得当，应力张量的守恒方程
\begin{equation}
  \tag{4.18}
  \label{eq:4-18}
  \partial_-T_{++}+\partial_+T_{+-}=0,\qquad \partial_+T_{--}+\partial_-T_{+-}=0
\end{equation}
（经典地，由于 $T_{+-}$ 为零，它们约化为 $\partial_-T_{++}=0$、$\partial_+T_{--}=0$）在量子力学上仍然成立。这样的守恒应力张量，如我们所知，导致庞加莱不变性。

让我们直接论证：在 $T_{+-}$ 的适当选取下，方程
\begin{equation}
  \tag{4.19}
  \label{eq:4-19}
  \partial_-T_{++}+\partial_+T_{+-}=0
\end{equation}
为何成立。为此回想：经典地，我们理论中的算符具有双分次（bigrading）$(d_+,d_-)$（关于 $x_+$、$x_-$ 的伸缩）。在量子理论中，由于共形不变性被破坏，沿反对角（antidiagonal）方向的次数差 $d_+-d_-$ 仍然保留，而总和 $d_++d_-$ 则变成一个滤子（filtration）。于是双分次变成了（双）滤子。

让我们对各种双次数的 $SO(N)$ 不变算符作分类（为顾及结构只是一个滤子这一事实，对给定的 $(d_+,d_-)$，我们将列出双次数为 $d_+-n,\ d_--n$、$n\ge0$ 的算符）：
\[
\begin{array}{ll}
(0,0): & 1\\
(0,1): & \text{无}\\
(1,0): & \text{无}\\
(2,0): & (\partial_+\phi)^2=T_{++}\\
(0,2): & (\partial_-\phi)^2=T_{--}\\
(1,1): & \partial_+\phi\,\partial_-\phi,\ 1\\
(2,1): & \partial_+^2\phi\,\partial_-\phi
\end{array}
\]

这张表意味着：在 $T_{ij}$ 的适当重整化下，(4.19) 必须成立，理由在于量纲。事实上，$\partial_-T_{++}$ 是双次数 $(2,1)$ 的算符，$\partial_+(\partial_+\phi\,\partial_-\phi)$ 也是，故这两个算符之间必存在一个线性关系。该线性关系并不是 $\partial_+(\partial_+\phi\,\partial_-\phi)=0$（因为这个算符经典地并不为零），所以该关系中 $\partial_-T_{++}$ 的系数非零，其形式即 (4.19)，其中 $T_{+-}$ 是 $\partial_+\phi\,\partial_-\phi$ 的某个倍数。

\textbf{注（Remark）.} 事实上可以证明 $T_{+-}=\beta\,\partial_+\phi\,\partial_-\phi$，其中 $\beta$ 是 $\beta$ 函数。于是，$\beta$ 函数度量了应力张量的迹，也度量了共形不变性被破坏的程度。

现在我们将用（量子的）应力-能量张量来构造一个被 Coleman–Mandula 定理禁戒的对称性。构造这一对称性的想法如下。设我们找到两个局域算符 $X_+,X_-$，使得
\[
\partial_+X_-=-\partial_-X_+,
\]
但并不存在 $Y$ 使 $X_-=-\partial_-Y$、$X_+=\partial_+Y$。此时，算符值的微分形式 $J=X_+dx_+-X_-dx_-$ 就是一个非平凡的守恒流（它不是任何 $Y$ 的 $dY$ 形式），而荷算符 $Q=\int_C J$（其中 $C$ 是类空闭链）守恒，并生成一个对称性。所以我们应当构造具有所述性质的 $X_+,X_-$。

% [Tutor Note: 表中 $(1,1)$ 一行同时列出 $\partial_+\phi\partial_-\phi$ 与 $1$，看起来奇怪，其实正是「滤子而非分次」的直接后果：可用的不变量只有 $d_+-d_-$，故一个朴素双次数为 $(1,1)$ 的算符可以与 $(0,0)$ 的单位算符混合（取 $n=1$），这正是 $\partial_+\phi\partial_-\phi$ 获得反常维数、并可与常数混合的体现。同理 $(2,1)$ 一行只列一个算符，是因为能与它混合的 $(1,0)$、$(0,-1)$ 等位置没有 $SO(N)$ 不变算符。若把滤子误当成分次（即认为双次数严格可加），就会得出「$\partial_-T_{++}$ 与 $\partial_+(\partial_+\phi\partial_-\phi)$ 无从相关」的错误结论，(4.19) 也就无从证明。]

% [Tutor Note: (4.18) 是 $\partial_\mu T^{\mu}{}_{\nu}=0$ 在光锥坐标下的两个分量，其中用到了 $T_{\mu\nu}$ 的对称性（故 $T_{-+}=T_{+-}$，第二式左端本应写作 $\partial_+T_{--}+\partial_-T_{-+}$）。作者只对第一式给出论证（即 (4.19)），第二式由 $x_+\leftrightarrow x_-$ 的交换对称性同样成立，这一步在原文中是默认的。核对指标时注意两式的结构必须完全平行：把第二式写成 $\partial_+T_{--}+\partial_-T_{-+}=0$ 之后，两个方程之间只差 $+\leftrightarrow-$ 的整体互换；若写成 $\partial_+T_{--}+\partial_+T_{+-}=0$ 或 $\partial_+T_{++}+\partial_-T_{+-}=0$，指标就没有配平。]
